% Bagian: turing-machines
% Bab: machines-computations
% Subbagian: representing-tms

\documentclass[../../../include/open-logic-section]{subfiles}

\begin{document}

\olfileid[id]{tur}{mac}{rep}
\olsection{Merepresentasikan Mesin Turing}

\begin{explain}
Mesin Turing dapat direpresentasikan secara visual dengan \emph{diagram
keadaan}. Diagram ini tersusun dari simpul-simpul keadaan yang dihubungkan oleh
panah. Tidak mengherankan bahwa setiap simpul keadaan merepresentasikan satu
keadaan mesin. Setiap panah merepresentasikan suatu instruksi yang dapat
dijalankan dari keadaan tersebut, dengan rincian instruksinya ditulis di atas
atau di bawah panah yang bersangkutan. Perhatikan mesin berikut, yang hanya
mempunyai dua keadaan internal, $q_0$ dan $q_1$, serta satu instruksi:
\[
\begin{tikzpicture}[->,>=stealth',shorten >=1pt,auto,node distance=2.8cm,
                    semithick]
  \tikzstyle{every state}=[fill=none,draw=black,text=black]

  \node[initial,state]   (A)              {$q_0$};
  \node[state]   (B) [right of=A] {$q_1$};

  \path (A) edge  node {\TMtrans{\TMblank}{\TMstroke}{\TMright}} (B);
\end{tikzpicture}
\]
Ingatlah bahwa mesin Turing mempunyai kepala baca-tulis dan pita yang memuat
masukan. Instruksi tersebut dapat dibaca sebagai \emph{jika sedang membaca
$\TMblank$ dalam keadaan~$q_0$, tulis~$\TMstroke$, bergerak ke kanan, dan
berpindah ke keadaan~$q_1$}. Ini setara dengan fungsi transisi yang memetakan
$\tuple{q_0, \TMblank}$ ke $\tuple{q_1, \TMstroke, \TMright}$.
\end{explain}

\begin{ex}
\emph{Mesin Genap}: Mesin Turing berikut berhenti jika dan hanya jika terdapat
sejumlah genap simbol~$\TMstroke$ pada pita (dengan asumsi bahwa semua simbol
$\TMstroke$ mendahului simbol~$\TMblank$ pertama pada pita).
\[
\begin{tikzpicture}[->,>=stealth',shorten >=1pt,auto,node distance=2.8cm,
                    semithick]
  \tikzstyle{every state}=[fill=none,draw=black,text=black]

  \node[initial,state]         (A)              {$q_0$};
  \node[state]         (B) [right of=A] {$q_1$};

  \path (A) edge [bend left] node {\TMtrans{\TMstroke}{\TMstroke}{\TMright}} (B)
        (B) edge [loop above] node {\TMtrans{\TMblank}{\TMblank}{\TMright}} (B)
            edge [bend left] node {\TMtrans{\TMstroke}{\TMstroke}{\TMright}} (A);
\end{tikzpicture}
\]

Diagram keadaan tersebut bersesuaian dengan fungsi transisi berikut:
\begin{align*}
\delta(q_0, \TMstroke) & = \tuple{q_1, \TMstroke, \TMright},\\
\delta(q_1, \TMstroke) & = \tuple{q_0, \TMstroke, \TMright},\\
\delta(q_1, \TMblank)  & = \tuple{q_1, \TMblank, \TMright}
\end{align*}
\end{ex}

\begin{explain}
Mesin di atas berhenti hanya apabila masukannya berupa sejumlah genap goresan.
Jika tidak, mesin itu (secara teoretis) terus beroperasi tanpa batas. Untuk
setiap mesin dan masukan, kita dapat menelusuri \emph{konfigurasi-konfigurasi}
mesin guna menentukan keluarannya. Kita akan memberikan definisi formal
konfigurasi nanti. Untuk sementara, secara intuitif konfigurasi dapat kita
pandang sebagai rangkaian diagram yang memperlihatkan keadaan mesin pada setiap
saat selama beroperasi. Konfigurasi memperlihatkan isi pita, keadaan mesin, dan
letak kepala baca-tulis.

Marilah kita telusuri konfigurasi mesin genap ketika dijalankan dengan masukan
empat simbol~$\TMstroke$. Dalam hal ini, kita mengharapkan mesin itu berhenti.
Kemudian kita akan menjalankan mesin dengan masukan tiga simbol~$\TMstroke$;
dalam hal ini mesin akan berjalan selamanya.

Mesin bermula dalam keadaan~$q_0$ sambil memindai simbol~$\TMstroke$ paling
kiri. Keadaan awal mesin dapat kita representasikan sebagai berikut:
\[
\TMendtape \TMstroke_0 \TMstroke \TMstroke \TMstroke \TMblank \ldots
\]
Konfigurasi di atas mudah dipahami. Seperti tampak, mesin bermula dalam
keadaan~$q_0$ sambil memindai simbol~$\TMstroke$ paling kiri. Hal ini
direpresentasikan dengan nama keadaan sebagai subskrip pada simbol
$\TMstroke$ pertama. Instruksi yang berlaku pada saat ini ialah
$\delta(q_0, \TMstroke) = \tuple{q_1, \TMstroke, \TMright}$, sehingga mesin
bergerak ke kanan pada pita dan berpindah ke keadaan~$q_1$.
\[
\TMendtape \TMstroke \TMstroke_1 \TMstroke \TMstroke \TMblank \ldots
\]
Karena mesin kini berada dalam keadaan~$q_1$ sambil memindai~$\TMstroke$, kita
harus ``mengikuti'' instruksi $\delta(q_1, \TMstroke) = \tuple{q_0,
  \TMstroke, \TMright}$. Hasilnya ialah konfigurasi
\[
\TMendtape \TMstroke \TMstroke \TMstroke_0 \TMstroke \TMblank \ldots
\]
Ketika mesin melanjutkan operasinya, aturan-aturan diterapkan lagi dengan
urutan yang sama dan menghasilkan dua konfigurasi berikut:
\[
\TMendtape \TMstroke \TMstroke \TMstroke \TMstroke_1 \TMblank \ldots
\]
\[
\TMendtape \TMstroke \TMstroke \TMstroke \TMstroke \TMblank_0 \ldots
\]
Mesin sekarang berada dalam keadaan~$q_0$ sambil memindai~$\TMblank$. Dari
diagram transisi, mudah dilihat bahwa tidak ada instruksi yang dapat dijalankan,
sehingga mesin telah berhenti. Ini berarti masukannya diterima.

Selanjutnya, misalkan mesin kita jalankan dengan masukan tiga
simbol~$\TMstroke$. Beberapa konfigurasi pertama serupa karena instruksi yang
sama dijalankan, hanya dengan sedikit perbedaan pada masukan pita:
\[
\TMendtape \TMstroke_0 \TMstroke \TMstroke \TMblank \ldots
\]
\[
\TMendtape \TMstroke \TMstroke_1 \TMstroke \TMblank \ldots
\]
\[
\TMendtape \TMstroke \TMstroke \TMstroke_0 \TMblank \ldots
\]
\[
\TMendtape \TMstroke \TMstroke \TMstroke \TMblank_1 \ldots
\]
Mesin kini telah melintasi semua simbol~$\TMstroke$ dan sedang membaca
$\TMblank$ dalam keadaan~$q_1$. Seperti ditunjukkan oleh diagram, terdapat
instruksi berbentuk $\delta(q_1, \TMblank) =\tuple{q_1,
\TMblank, \TMright}$. Karena pita terisi simbol~$\TMblank$ tanpa batas ke
kanan, mesin akan terus menjalankan instruksi ini \emph{selamanya}, tetap dalam
keadaan~$q_1$ dan bergerak semakin jauh ke kanan. Mesin tidak akan pernah
berhenti dan tidak menerima masukan tersebut.
\end{explain}

\begin{explain}
Penting diperhatikan bahwa tidak semua mesin akan berhenti. Jika berhenti
berarti mesin kehabisan instruksi untuk dijalankan, kita dapat membuat mesin
yang tidak pernah berhenti cukup dengan memastikan bahwa terdapat panah keluar
bagi setiap simbol pada setiap keadaan. Mesin genap dapat diubah agar berjalan
tanpa batas dengan menambahkan instruksi untuk memindai~$\TMblank$ pada~$q_0$.
\end{explain}

\begin{ex}
\[
\begin{tikzpicture}[->,>=stealth',shorten >=1pt,auto,node distance=2.8cm,
                    semithick]
  \tikzstyle{every state}=[fill=none,draw=black,text=black]

  \node[initial,state] (A)              {$q_0$};
  \node[state]         (B) [right of=A] {$q_1$};

  \path
  (A) edge [bend left] node {\TMtrans{\TMstroke}{\TMstroke}{\TMright}} (B)
      edge [loop above] node {\TMtrans{\TMblank}{\TMblank}{\TMright}} (A)
  (B) edge [loop above] node {\TMtrans{\TMblank}{\TMblank}{\TMright}} (B)
      edge [bend left] node {\TMtrans{\TMstroke}{\TMstroke}{\TMright}} (A);
\end{tikzpicture}
\]
\end{ex}

\begin{explain}
Tabel mesin merupakan cara lain untuk merepresentasikan mesin Turing. Dalam
tabel mesin, alfabet pita ditampilkan pada sumbu~$x$ dan himpunan keadaan mesin
pada sumbu~$y$. Di dalam tabel, pada perpotongan setiap keadaan dan simbol,
dituliskan bagian lain dari instruksi---keadaan baru, simbol baru, dan arah
gerak. Tabel mesin memudahkan kita menentukan pada keadaan apa, dan untuk
simbol apa, mesin berhenti. Setiap celah dalam tabel merupakan titik yang
mungkin membuat mesin berhenti. Berbeda dari diagram keadaan dan himpunan
instruksi, yang titik-titik penghentiannya tidak selalu langsung tampak,
setiap titik penghentian dapat segera dikenali dengan mencari celah dalam
tabel mesin.
\end{explain}

\begin{ex}
Tabel mesin bagi mesin genap adalah:
\[
\centering
\begin{tabular}{lllll}
\cline{2-4}
\multicolumn{1}{l|}{}      & \multicolumn{1}{c|}{$\TMblank$}
& \multicolumn{1}{c|}{$\TMstroke$}     & \multicolumn{1}{c|}{$\TMendtape$}   &  \\ \cline{1-4}
\multicolumn{1}{|l|}{$q_0$} & \multicolumn{1}{l|}{}
& \multicolumn{1}{l|}{\TMtrans{\TMstroke}{q_1}{\TMright}} &
\multicolumn{1}{l|}{\phantom{\TMtrans{\TMstroke}{q_1}{\TMright}}} &  \\ \cline{1-4}
\multicolumn{1}{|l|}{$q_1$} & \multicolumn{1}{l|}{\TMtrans{\TMblank}{q_1}{\TMright}}
& \multicolumn{1}{l|}{\TMtrans{\TMstroke}{q_0}{\TMright}} &
\multicolumn{1}{l|}{\phantom{\TMtrans{\TMstroke}{q_1}{\TMright}}} &  \\ \cline{1-4}
\end{tabular}
\]
Seperti dapat kita lihat, mesin berhenti ketika memindai~$\TMblank$ dalam
keadaan~$q_0$.
\end{ex}

\begin{explain}
Sejauh ini kita hanya membahas mesin yang membaca dan menerima masukan. Namun,
mesin Turing mampu membaca sekaligus menulis. Salah satu contohnya (walaupun
terdapat banyak sekali contoh) ialah \emph{mesin pengganda}. Ketika dijalankan
dengan sebuah blok yang terdiri atas~$n$ simbol~$\TMstroke$ pada pita, mesin
pengganda menghasilkan sebuah blok yang terdiri atas~$2n$ simbol~$\TMstroke$.
\end{explain}

\begin{ex}
  \ollabel{ex:doubler} Sebelum membangun mesin pengganda, penting untuk
  menyusun suatu \emph{strategi} pemecahan masalah. Karena mesin (sebagaimana
  telah kita rumuskan) tidak dapat mengingat berapa banyak simbol~$\TMstroke$
  yang telah dibacanya, kita perlu menemukan cara melacak semua simbol
  $\TMstroke$ pada pita. Salah satu caranya ialah memisahkan keluaran dari
  masukan dengan sebuah~$\TMblank$. Mesin kemudian dapat menghapus simbol
  $\TMstroke$ pertama dari masukan, melintasi sisa masukan, meninggalkan sebuah
  $\TMblank$, lalu menulis dua simbol~$\TMstroke$ baru. Selanjutnya mesin
  kembali dan mencari simbol~$\TMstroke$ kedua dalam masukan, lalu
  menggandakannya juga. Untuk setiap satu simbol~$\TMstroke$ pada masukan,
  mesin menulis dua simbol~$\TMstroke$ pada keluaran. Dengan menghapus masukan
  sedikit demi sedikit selama operasi, kita dapat memastikan bahwa tidak ada
  simbol~$\TMstroke$ yang terlewat atau digandakan dua kali. Setelah seluruh
  masukan dihapus, akan tersisa $2n$ simbol~$\TMstroke$ pada pita. Diagram
  keadaan mesin Turing yang dihasilkan ditampilkan pada \olref{fig:doubler}.
  \begin{figure}
\[
\begin{tikzpicture}[->,>=stealth',shorten >=1pt,auto,node distance=2.8cm,
                    semithick]
  \tikzstyle{every state}=[fill=none,draw=black,text=black]

  \node[initial,state] (1)              {$q_0$};
  \node[state]         (2) [right of=1] {$q_1$};
  \node[state]         (3) [right of=2] {$q_2$};
  \node[state]         (4) [below of=3] {$q_3$};
  \node[state]         (5) [left of=4]  {$q_4$};
  \node[state]         (6) [left of=5]  {$q_5$};

  \path (1) edge node {\TMtrans{\TMstroke}{\TMblank}{\TMright}} (2)
    (2) edge [loop above] node {\TMtrans{\TMstroke}{\TMstroke}{\TMright}} (2)
      edge node {\TMtrans{\TMblank}{\TMblank}{\TMright}} (3)
    (3) edge [loop above] node {\TMtrans{\TMstroke}{\TMstroke}{\TMright}} (3)
        edge node {\TMtrans{\TMblank}{\TMstroke}{\TMright}} (4)
    (4) edge [loop below] node {\TMtrans{\TMblank}{\TMstroke}{\TMleft}} (4)
        edge node {\TMtrans{\TMstroke}{\TMstroke}{\TMleft}} (5)
    (5) edge [loop below]  node {\TMtrans{\TMstroke}{\TMstroke}{\TMleft}} (5)
        edge              node {\TMtrans{\TMblank}{\TMblank}{\TMleft}} (6)
    (6) edge [loop below] node {\TMtrans{\TMstroke}{\TMstroke}{\TMleft}} (6)
        edge              node {\TMtrans{\TMblank}{\TMblank}{\TMright}} (1);
\end{tikzpicture}
\]
\caption{Mesin pengganda}
\ollabel{fig:doubler}
\end{figure}
\end{ex}

\begin{prob}
Pilih suatu masukan sebarang dan telusuri konfigurasi-konfigurasi mesin
pengganda pada \olref[tur][mac][rep]{ex:doubler}.
\end{prob}

\begin{prob}
Rancang mesin Turing dengan alfabet $\{\TMendtape,\TMblank, A, B\}$ yang
menerima, yakni berhenti pada, setiap untai $A$ dan $B$ yang jumlah simbol
$A$-nya sama dengan jumlah simbol $B$-nya \emph{dan} semua $A$ mendahului
semua $B$; serta menolak, yakni tidak berhenti pada, setiap untai yang jumlah
$A$-nya tidak sama dengan jumlah $B$-nya atau semua $A$ tidak mendahului semua
$B$. (Misalnya, mesin harus menerima $AABB$ dan $AAABBB$, tetapi menolak
$AAB$ maupun $AABBAABB$.)
\end{prob}

\begin{prob}
Rancang mesin Turing dengan alfabet $\{\TMendtape,\TMblank, A, B\}$ yang
menerima sebarang untai~$\alpha$ yang terdiri atas $A$ dan $B$ sebagai masukan,
lalu menduplikasinya sehingga menghasilkan keluaran berbentuk
$\alpha\alpha$. (Misalnya, masukan $ABBA$ harus menghasilkan keluaran
$ABBAABBA$.)
\end{prob}

\begin{prob}
\emph{Terurut alfabetis?:} Rancang mesin Turing dengan alfabet
$\{\TMendtape,\TMblank, A, B\}$ yang, ketika diberi masukan berupa barisan
berhingga $A$ dan $B$, memeriksa apakah semua $A$ muncul di sebelah kiri semua
$B$. Mesin harus membiarkan untai masukan tetap pada pita, lalu berhenti jika
untai itu ``terurut alfabetis'' atau berulang selamanya jika tidak.
\end{prob}

\begin{prob}
\emph{Pengurut alfabetis:} Rancang mesin Turing dengan alfabet
$\{\TMendtape,\TMblank, A, B\}$ yang menerima masukan berupa barisan berhingga
$A$ dan $B$, lalu menyusun ulang simbol-simbol tersebut sehingga semua $A$
berada di sebelah kiri semua $B$. (Misalnya, barisan $BABAA$ harus menjadi
$AAABB$, dan barisan $ABBABB$ harus menjadi $AABBBB$.)
\end{prob}

\end{document}
